
L'approche BFS génère deux informations distinctes et nécessaires : la distance minimale (un naturel) qui valide l'existence d'un chemin, et la trace des prédécesseurs (un tableau) qui permet de remonter de l'arrivée vers le départ pour reconstituer l'itinéraire. Le langage C ne permettant pas le retour de valeurs multiples, l'encapsulation de ces deux données dans la structure ResultatBFS permet de récupérer l'intégralité du résultat de l'exploration en une seule opération.

Nous avons donc ajouté une structure afin de nous aider dans notre démarche.

\item[$\bullet$] 
        \begin{enregistrement}{ResultatBFS}
            \champEnregistrement{Distance}{\naturel}             
            \champEnregistrement{predecesseurs}{\tableauUneDimension{ }de {\naturelNonNul}    }      
        \end{enregistrement}
        
\vspace{1cm}

\begin{algorithme}
    \fonction{genererCheminOptimal}
        {v : Ville}
        {Chemin}
        {nb\_cases, depart, nb\_objectifs, i ,j , k ,l, taille\_max , idx\_depart , idx\_arrivee : \naturelNonNul \\
        longueur\_totale, longueur\_segment, debut  : \naturel\\
        min\_cout  : \entier\\
        distances : \tableauDeuxDimensions{ } { } {de} {\entier} \\
        objectifs , sequence\_temp, meilleure\_sequence, chemin\_final, segment : \tableauUneDimension{ } de {\naturelNonNul} \\
        visites : \tableauUneDimension{ } de {\bool}\\
        predecesseurs : \tableauDeuxDimensions{ } { } {de} {\naturel} \\
        resBFS : ResultatBFS \\
        resultat : Chemin \\
        } 
        {
        \sialors{v = \motclefPointeurNULL}
            {\retourner{creerChemin()}}

        \affecter{nb\_cases}{obtenirNbCases(v)}
        \affecter{depart}{obtenirNumCaseDepart(v)}

        \affecter{objectifs}{allouer(nb\_cases + 1)}
        \affecter{objectifs[0]}{depart}
        \affecter{nb\_objectifs}{1}

        \pour{i}{1}{nb\_cases}{}{
            \affecter{c}{obtenirCase(v, i)}
            \sialors{estObligatoire(c) et i $\neq$ depart}
            {
                \affecter{objectifs[nb\_objectifs]}{i}
                \affecter{nb\_objectifs}{nb\_objectifs + 1}
            }
        }

        \sialors{nb\_objectifs $\leq$ 1}
        {
            \affecter{resultat}{construire\_chemin\_final(v, depart, 0, depart, obtenirDirDepart(v))}
            \instruction{desallouer(objectifs)}
            \retourner{resultat}
        }

        \affecter{distances}{allouerMatrice(nb\_objectifs, nb\_objectifs)}
        \affecter{predecesseurs}{allouerMatrice(nb\_objectifs, nb\_objectifs)}

        \pour{i}{0}{nb\_objectifs - 1}{}{
            \affecter{distances[i]}{allouer(nb\_objectifs)}
            \affecter{predecesseurs[i]}{allouer(nb\_objectifs)}

            \pour{j}{0}{nb\_objectifs - 1}{}{
                \sialorssinon{i = j}
                {
                    \affecter{distances[i][j]}{0}
                    \affecter{predecesseurs[i][j]}{\motclefPointeurNULL}
                }
                {
                    \affecter{resBFS}{bfs\_chemin\_court(v, objectifs[i], objectifs[j])}
                    \affecter{distances[i][j]}{resBFS.distance}
                    \affecter{predecesseurs[i][j]}{resBFS.predecesseurs}

                    \sialors{resBFS.distance = {\infty }}
                    {
                        \instruction{desallouer(objectifs)}
                        \pour{k}{0}{i}{}{
                            \pour{l}{0}{nb\_objectifs - 1}{}{
                                \sialors{predecesseurs[k][l] $\neq$ \motclefPointeurNULL}
                                    {\instruction{desallouer(predecesseurs[k][l])}}
                            }
                            \instruction{desallouer(predecesseurs[k])}
                            \instruction{desallouer(distances[k])}
                        }
                        \instruction{desallouer(predecesseurs)}
                        \instruction{desallouer(distances)}
                        \retourner{creerChemin()}
                    }
                }
            }
        }

        \affecter{sequence\_temp}{allouer(nb\_objectifs)}
        \affecter{meilleure\_sequence}{allouer(nb\_objectifs)}
        \affecter{visites}{allouer(nb\_objectifs)}
        \affecter{min\_cout}{\infty}

        \affecter{sequence\_temp[0]}{0}
        \affecter{visites[0]}{vrai}

        \instruction{explorer\_ordre\_objectifs(1, nb\_objectifs,
                     sequence\_temp, visites, 0,
                     distances, meilleure\_sequence, \&min\_cout)}

        \sialors{min\_cout = \infty}
        {
            \instruction{desallouer(sequence\_temp)}
            \instruction{desallouer(meilleure\_sequence)}
            \instruction{desallouer(visites)}
            \instruction{desallouer(objectifs)}
            \pour{i}{0}{nb\_objectifs - 1}{}{
                \pour{j}{0}{nb\_objectifs - 1}{}{
                    \sialors{predecesseurs[i][j] $\neq$ \motclefPointeurNULL}
                        {\instruction{desallouer(predecesseurs[i][j])}}
                }
                \instruction{desallouer(predecesseurs[i])}
                \instruction{desallouer(distances[i])}
            }
            \instruction{desallouer(predecesseurs)}
            \instruction{desallouer(distances)}
            \retourner{creerChemin()}
        }

        \affecter{taille\_max}{2 \times {nb\_cases} + 2}
        \affecter{chemin\_final}{allouer(taille\_max)}
        \affecter{longueur\_totale}{0}

        \pour{i}{0}{nb\_objectifs - 1}{}{
            \affecter{idx\_depart}{meilleure\_sequence[i]}
            \sialorssinon{i = nb\_objectifs - 1}
                {\affecter{idx\_arrivee}{0}}
                {\affecter{idx\_arrivee}{meilleure\_sequence[i + 1]}}

            \affecter{segment}{allouer(nb\_cases + 1)}
            \affecter{longueur\_segment}{reconstruire\_chemin(
                objectifs[idx\_depart],
                objectifs[idx\_arrivee],
                predecesseurs[idx\_depart][idx\_arrivee],
                segment)}

            \sialors{longueur\_segment < 0}
            {
                \instruction{desallouer(segment)}
                \continuer
            }

            \sialorssinon{i = 0}
                {\affecter{debut}{0}}
                {\affecter{debut}{1}}

            \pour{j}{debut}{longueur\_segment}{}{
                \affecter{chemin\_final[longueur\_totale]}{segment[j]}
                \affecter{longueur\_totale}{longueur\_totale + 1}
            }

            \instruction{desallouer(segment)}
        }

        \affecter{resultat}{construire\_chemin\_final(
            v, chemin\_final, longueur\_totale - 1,
            depart, obtenirDirDepart(v))}

        \instruction{desallouer(objectifs)}
        \instruction{desallouer(chemin\_final)}
        \instruction{desallouer(sequence\_temp)}
        \instruction{desallouer(meilleure\_sequence)}
        \instruction{desallouer(visites)}

        \pour{i}{0}{nb\_objectifs - 1}{}{
            \pour{j}{0}{nb\_objectifs - 1}{}{
                \sialors{i $\neq$ j et predecesseurs[i][j] $\neq$ \motclefPointeurNULL}
                    {\instruction{desallouer(predecesseurs[i][j])}}
            }
            \instruction{desallouer(predecesseurs[i])}
            \instruction{desallouer(distances[i])}
        }

        \instruction{desallouer(predecesseurs)}
        \instruction{desallouer(distances)}

        \retourner{resultat}
    }
\end{algorithme}




